\section{Complete experimental and implementation protocol}
\label{app:experimental-protocol}

This appendix records the protocol that turns the executable-process
definitions in the main text into claim-bearing experiments.  It is deliberately
more specific than a conventional implementation summary.  A legal executor
can make validity nearly automatic while still producing an uninformative
transition law; conversely, an optimizer can improve endpoints while obscuring
an unusable base process.  We therefore separate support, learning, control,
and task-level evidence throughout.

\subsection{Evidence, freezing, and attempted-sample accounting}
\label{app:evidence-freeze}

\paragraph{Evidence classes.}
Each reported item is classified as a mathematical derivation, a frozen
implementation fact, a registered protocol, or a measured result.  Result
artifacts are accepted only when they contain the expected schema, physical
file hashes, implementation and configuration identities, partition identity,
seed, and exact checkpoint hash.  A development artifact may validate code but
cannot fill a final result cell.  Historical measurements are labeled
diagnostic and do not select the production checkpoint.

\paragraph{Partition use.}
Training data fit parameters.  Development panels are used to repair
factorizations and set pilot thresholds.  Validation data choose checkpoints
and controller hyperparameters under the frozen rules below.  The final test
partition is opened once after the checkpoint, tasks, thresholds, and
comparators are frozen.  Because diagnostics from the completed 16,000-step
run have already been inspected, a repaired model uses a new sealed final
holdout.

\paragraph{Denominators.}
Every attempted trajectory remains in the denominator for validity, target
success, constraint satisfaction, and budget compliance.  A failed proposal,
zero-value control state, exhausted successor fiber, timeout, or external
baseline failure is recorded under a distinct reason code rather than silently
resampled.  Molecules that fail independent standardization or evaluation are
reported separately and count as failed task attempts whenever that evaluator
is part of the registered endpoint.

\paragraph{Compute.}
The primary control cost is the number of property-oracle calls.  We also
record model forward passes, raw marks scored, canonical successors scored,
particles propagated, accepted productive edits, executor calls, wall-clock
time, accelerator hours, and peak memory.  Training reports optimizer steps,
examples, unique exact states, successor-cache hits, and effective
family/capability coefficients.  A controller that evaluates $m$ candidates or
propagates $m$ particles pays for all $m$, not only for the committed
successor.

\subsection{Editing corpus construction}
\label{app:editing-corpus}

\paragraph{Training target.}
The editing model is a property-agnostic prior over canonical molecular
successors.  Source molecule, goal, remaining budget, protected mapping, and
path-constraint state are controller inputs and are not leaked into the base
prior.  Each stored trace contains exact persistent-slot states and executable
actions.  Canonical strings are molecular identity metadata; training never
reconstructs a slot-addressed state by reparsing a canonical string.

\begin{table}[htbp]
  \centering
  \caption{\textbf{Editing-corpus evidence lanes.} ``Real endpoint'' means both
  endpoints occur in the underlying molecular source; intermediate states are
  still executor-certified.  Final counts and effective coefficients are
  filled only from the frozen V2 inventory.}
  \label{tab:appendix-corpus-lanes}
  \small
  \begin{tabularx}{\linewidth}{@{}lYYY@{}}
    \toprule
    Lane & Purpose & Required evidence & Frozen inventory \\
    \midrule
    Observed local analogue &
    substitutions, bioisosteres, and local size changes &
    real endpoint; shared-source branches where available &
    \resultpending{APP_CORPUS_LOCAL} \\
    Operator-aware real endpoint &
    prefer direct restate, reorder, reroute, birth/death, and cycle programs
    over delete--rebuild fallbacks &
    executable direct path or explicit compiler failure reason &
    \resultpending{APP_CORPUS_OPERATOR_AWARE} \\
    Linker/positional/topology analogue &
    attachment relocation, linker changes, and topology changes &
    real endpoint preferred; certified synthetic supplement allowed &
    \resultpending{APP_CORPUS_LINKER_TOPOLOGY} \\
    Observed series path &
    short paths, shared prefixes, and alternative routes through observed
    analogue series &
    molecule-, scaffold-, series-, and source-group provenance &
    \resultpending{APP_CORPUS_SERIES} \\
    Reversible synthetic walk &
    rare contexts, inverse/correction coverage, and mixed-family composition &
    production executor, closure, charge policy, and exact replay &
    \resultpending{APP_CORPUS_REVERSIBLE} \\
    \bottomrule
  \end{tabularx}
\end{table}

The lanes are evidence categories rather than automatically equal mixture
components.  In particular, a forward/reverse encoding of one matched pair is
not by itself an observed series, a branch group, or a dynamic correction
example.  Conversely, the local Markov prior need not memorize complete
Pareto-optimal trajectories: controllers compose local transitions at
inference.  The final corpus census therefore determines which multi-step
lanes add measured coverage beyond real local endpoints and reversible
executor walks; an aspirational field is not allowed to block a bounded pilot
without a demonstrated capability gap.

\paragraph{Compiler order.}
For a real endpoint pair, the compiler first searches direct semantic
operations supported by the production factorization.  Only then may it use a
multi-step delete/insert or other primitive composition.  Both directions are
compiled independently because executability and path length can differ.
Every retained trace replays from its stored exact source to its stored exact
target.  Unsupported pairs remain in an audit inventory with a typed failure
reason.

\paragraph{Splits.}
Molecule identity, scaffold component, series group, upstream
document/source group, and transformation signature are group-level split
keys.  Both endpoints of a pair must belong to the same partition;
cross-partition pairs are dropped before trace generation.  The frozen
manifest reports molecule, scaffold, series, and source-group overlaps both
within each layer and across layers.  Manifest identity is computed from the
bytes of every packed shard, not only its path and row count.

\paragraph{Sampling law.}
Sampling is hierarchical:
\[
  \text{data lane}
  \rightarrow
  \text{series/scaffold/source group}
  \rightarrow
  \text{semantic capability cell}
  \rightarrow
  \text{endpoint program}
  \rightarrow
  \text{progress state}.
\]
Raw-pair-uniform sampling is disallowed because common delete/insert traces
would dominate rare but load-bearing cycle and reroute events.  For every run
we reconstruct the realized coefficient of each teacher family, capability
cell, and path position.  Packed transition addressability is an
implementation optimization and does not change this hierarchy.

\begin{table}[htbp]
  \centering
  \caption{\textbf{Frozen editing-corpus census.} Counts are reported after all
  representability, charge, split, and whole-trace support exclusions.}
  \label{tab:appendix-corpus-census}
  \small
  \begin{tabularx}{\linewidth}{@{}lYYYY@{}}
    \toprule
    Layer/lane & Traces & Exact states & Unique sources & Unique scaffolds \\
    \midrule
    General corruption &
    \resultpending{APP_CENSUS_CORRUPTION_TRACES} &
    \resultpending{APP_CENSUS_CORRUPTION_STATES} &
    \resultpending{APP_CENSUS_CORRUPTION_SOURCES} &
    \resultpending{APP_CENSUS_CORRUPTION_SCAFFOLDS} \\
    Cycle operations &
    \resultpending{APP_CENSUS_CYCLE_TRACES} &
    \resultpending{APP_CENSUS_CYCLE_STATES} &
    \resultpending{APP_CENSUS_CYCLE_SOURCES} &
    \resultpending{APP_CENSUS_CYCLE_SCAFFOLDS} \\
    Real-endpoint analogues &
    \resultpending{APP_CENSUS_ANALOGUE_TRACES} &
    \resultpending{APP_CENSUS_ANALOGUE_STATES} &
    \resultpending{APP_CENSUS_ANALOGUE_SOURCES} &
    \resultpending{APP_CENSUS_ANALOGUE_SCAFFOLDS} \\
    Reversible/mixed-family supplements &
    \resultpending{APP_CENSUS_SUPPLEMENT_TRACES} &
    \resultpending{APP_CENSUS_SUPPLEMENT_STATES} &
    \resultpending{APP_CENSUS_SUPPLEMENT_SOURCES} &
    \resultpending{APP_CENSUS_SUPPLEMENT_SCAFFOLDS} \\
    \bottomrule
  \end{tabularx}
\end{table}

\subsection{Production support and chemistry audits}
\label{app:chemistry-audits}

\paragraph{State contract.}
The molecular state uses at most 40 persistent slots and the declared
broad-organic element--valence vocabulary.  Active atoms are identified by the
authoritative element predicate; neither nonzero atom-state indices nor a
prefix slice of the padded tensor defines the molecule.  The SCAR bookkeeping
state is occupied but is not an element.  Graph cycle rank is always
$|E|-|V|+c$ and is reported separately from SSSR ring count and ring-system
count.

\paragraph{Charge policy.}
Formal-charge design is outside the current action space.  A legal editing
transition preserves the complete persistent-slot charge array and protects
the element, implicit hydrogen count, and full bond row of every initially
charged slot.  The same predicate must filter training traces, the model's
legal candidate masks, successor-cache construction, sampling, and evaluation.
Protecting only net charge is insufficient.  The final support audit enumerates
charged-source candidates and requires zero policy-violating probability mass:
\resultpending{APP_CHARGE_SUPPORT_AUDIT}.

\paragraph{Operator basis.}
The bounded rebuild tests the eight families in
Table~\ref{tab:operators}.  Ring-system deletion and whole-ring growth are
disabled.  Ring-system restatement is a coordinated accelerator under the
current primitive bond-order masks; its inclusion in the bounded pilot does
not authorize it as final production support.  Final support is frozen only
after successor-level capacity, retention, rollout, and topology-task audits.

\paragraph{Birth and death.}
Insertion is either a root birth from the null state or a one-neighbor
connected birth.  General simultaneous multi-neighbor insertion is absent.
Deleting a degree-two atom may therefore lack a one-step insertion inverse,
although its endpoint can sometimes be reached by birth followed by cycle
closure and restatement.  We report primitive path cost instead of claiming
one-step inverse closure.

\paragraph{Cycles.}
Cycle closure adds one executable bond and cycle opening deletes one executable
ring edge.  Whole-ring templates do not define support.  Candidate audits
report teacher presence, exact execution, undirected endpoint
canonicalization, raw mark count, distinct canonical successor count, alias
multiplicity, and candidate-count-matched random baselines separately for
closure and opening.

\subsection{No-waste training and checkpoint selection}
\label{app:training-gates}

Long optimization is authorized only after the staged gates in
Table~\ref{tab:appendix-training-gates}.  A gate failure is diagnostic evidence,
not a reason to relax its threshold after viewing the result.

\begin{table}[htbp]
  \centering
  \caption{\textbf{Training release ladder.} A later stage cannot compensate
  for a failed earlier invariant.}
  \label{tab:appendix-training-gates}
  \small
  \begin{tabularx}{\linewidth}{@{}lYYY@{}}
    \toprule
    Stage & Purpose & Required evidence & Failure action \\
    \midrule
    Gate 0 &
    deterministic construction and support parity &
    exact corpus/cache/config hashes; teacher recovery; legal masks; nonempty
    active families &
    repair data or implementation \\
    T1 &
    bounded successor-level capacity &
    tiny unique-state and repeated-state panels; family NLL, probability,
    rank, candidates, aliases &
    repair labels/factorization/head \\
    P50 &
    gradient and collapse sentinel &
    50 exact updates; every family exposed; finite gradients; no probability
    collapse; immutable completion &
    repair recipe or sampling \\
    P500 &
    early learning and retention &
    successor improvement in every load-bearing family; inherited-capability
    sentinel when applicable &
    rebalance, distill, or change learning rates \\
    P2000 &
    generalization, rollouts, and calibration &
    held-source/scaffold/topology metrics; productive rollouts; no divergence &
    stop before long run \\
    Full run &
    final model fitting &
    all prior gates; frozen stopping/selection rule; resumable scientific
    identity &
    abort or keep preregistered best checkpoint \\
    \bottomrule
  \end{tabularx}
\end{table}

\paragraph{T1 panels.}
Panel construction is whole-trace support aware.  Filtering one displayed
teacher row is not sufficient when another progress state in the same trace
uses a disabled family: every progress transition needed by cache
materialization must belong to the exact active support.  Unique-state panels
test memorization of unambiguous local targets.  Repeated-state panels compare
the learned successor distribution with the empirical teacher distribution
and its entropy floor.

\paragraph{Long-run monitoring.}
Each recovery checkpoint contains model and optimizer state, scheduler
position, random-number states, completed optimizer steps, corpus and shard
inventory hashes, representability and charge overlays, successor-cache
identity, operator freeze, objective and sampler versions, sentinel panel
identity, and the frozen threshold/stopping policy.  A resume with any
different scientific field fails closed.

At every registered checkpoint we compute production-weighted and
family-balanced canonical-successor NLL; per-family NLL, probability, rank,
candidate count, successor count, and alias multiplicity; required-capability
retention; productive rollout statistics; and, where applicable, hazard
calibration.  The monitor returns one of three statuses:
\emph{continue}, \emph{stop and retain the registered best validation
checkpoint}, or \emph{abort for invariant/capability failure}.  Test metrics
never influence these decisions.

\begin{table}[htbp]
  \centering
  \caption{\textbf{Operator-resolved validation report for the selected editing
  checkpoint.} Mark accuracy is auxiliary; successor NLL and probability are
  the primary local-capacity measures.}
  \label{tab:appendix-family-diagnostics}
  \small
  \begin{tabularx}{\linewidth}{@{}lYYYYY@{}}
    \toprule
    Family & Succ. NLL $\downarrow$ & Teacher mass $\uparrow$ & MRR $\uparrow$ &
    Raw/canonical candidates & Alias multiplicity \\
    \midrule
    Atom insert & \resultpending{APP_FAM_INSERT_NLL} &
    \resultpending{APP_FAM_INSERT_PROB} & \resultpending{APP_FAM_INSERT_MRR} &
    \resultpending{APP_FAM_INSERT_CAND} & \resultpending{APP_FAM_INSERT_ALIAS} \\
    Atom delete & \resultpending{APP_FAM_DELETE_NLL} &
    \resultpending{APP_FAM_DELETE_PROB} & \resultpending{APP_FAM_DELETE_MRR} &
    \resultpending{APP_FAM_DELETE_CAND} & \resultpending{APP_FAM_DELETE_ALIAS} \\
    Atom restate & \resultpending{APP_FAM_RESTATE_NLL} &
    \resultpending{APP_FAM_RESTATE_PROB} & \resultpending{APP_FAM_RESTATE_MRR} &
    \resultpending{APP_FAM_RESTATE_CAND} & \resultpending{APP_FAM_RESTATE_ALIAS} \\
    Bond reorder & \resultpending{APP_FAM_REORDER_NLL} &
    \resultpending{APP_FAM_REORDER_PROB} & \resultpending{APP_FAM_REORDER_MRR} &
    \resultpending{APP_FAM_REORDER_CAND} & \resultpending{APP_FAM_REORDER_ALIAS} \\
    Bond reroute & \resultpending{APP_FAM_REROUTE_NLL} &
    \resultpending{APP_FAM_REROUTE_PROB} & \resultpending{APP_FAM_REROUTE_MRR} &
    \resultpending{APP_FAM_REROUTE_CAND} & \resultpending{APP_FAM_REROUTE_ALIAS} \\
    Cycle close & \resultpending{APP_FAM_CLOSE_NLL} &
    \resultpending{APP_FAM_CLOSE_PROB} & \resultpending{APP_FAM_CLOSE_MRR} &
    \resultpending{APP_FAM_CLOSE_CAND} & \resultpending{APP_FAM_CLOSE_ALIAS} \\
    Cycle open & \resultpending{APP_FAM_OPEN_NLL} &
    \resultpending{APP_FAM_OPEN_PROB} & \resultpending{APP_FAM_OPEN_MRR} &
    \resultpending{APP_FAM_OPEN_CAND} & \resultpending{APP_FAM_OPEN_ALIAS} \\
    Ring restate & \resultpending{APP_FAM_RING_RESTATE_NLL} &
    \resultpending{APP_FAM_RING_RESTATE_PROB} &
    \resultpending{APP_FAM_RING_RESTATE_MRR} &
    \resultpending{APP_FAM_RING_RESTATE_CAND} &
    \resultpending{APP_FAM_RING_RESTATE_ALIAS} \\
    \bottomrule
  \end{tabularx}
\end{table}

\paragraph{Checkpoint selection.}
All registered checkpoints are evaluated on the fixed validation corpus.
Primary selection minimizes production-weighted canonical-successor NLL;
balanced-family successor NLL is secondary.  Held-source, scaffold, topology,
charge-support, calibration, and productive-rollout safeguards can veto a
candidate.  Hazard metrics cannot select the fixed-budget editing checkpoint.
The complete leaderboard is frozen before the final test is opened.

\subsection{E2--E4: transport and structural adaptation}
\label{app:e2-e4-protocols}

\paragraph{E2 matched-law comparison.}
The learned and uniform arms use identical sources, canonical successor sets,
executor, random seeds, and productive-edit budgets.  The uniform arm assigns
$1/|\mathcal N(x)|$ to distinct molecular successors, not to marks.  Target
recovery credits a trajectory when its endpoint canonical identity equals the
held-out target within the registered budget.  Path overhead compares accepted
productive edits with the frozen verified compiler path or, on the exact
subset, the shortest path.  We separately report immediate reversals, repeated
states, longer cycles, net and gross atom/bond displacement, and termination
due to an empty productive fiber.

\paragraph{E3 cardinality.}
Sources and target atom-count intervals are stratified by initial count,
requested signed displacement, scaffold, and whether the target is impossible
under fixed cardinality.  Full \method is compared with no insertion, no
deletion, and fixed-cardinality support.  Dynamic trials first target a larger
interval and then switch to a smaller interval at the registered step.  Every
intermediate count, first hitting time, overshoot, correction count, source
similarity, and oracle call is retained.

\paragraph{E4 topology.}
Tasks include cycle creation, opening, ring-size change, topology
simplification, fused-system modification, and a create-then-reverse switch.
Spiro and bridged strata are included only when an independent reachability
builder finds a legal path under the frozen operator set.  Results are
stratified by graph cycle-rank delta, ring-size transition, ring-system count,
and fused/spiro/bridged category.  The no-cycle ablation removes only cycle
close/open; it does not reintroduce a ring template macro.

\begin{table}[htbp]
  \centering
  \caption{\textbf{Detailed size and topology adaptation.} The main text
  reports aggregates; this table preserves the registered difficult strata.}
  \label{tab:appendix-adaptation-strata}
  \small
  \begin{tabularx}{\linewidth}{@{}lYYYY@{}}
    \toprule
    Stratum & Sources & Target success $\uparrow$ & Valid path $\uparrow$ &
    Edits/oracles $\downarrow$ \\
    \midrule
    Grow & \resultpending{APP_E3_GROW_N} &
    \resultpending{APP_E3_GROW_SUCCESS} &
    \resultpending{APP_E3_GROW_VALID} & \resultpending{APP_E3_GROW_COST} \\
    Shrink & \resultpending{APP_E3_SHRINK_N} &
    \resultpending{APP_E3_SHRINK_SUCCESS} &
    \resultpending{APP_E3_SHRINK_VALID} & \resultpending{APP_E3_SHRINK_COST} \\
    Grow then shrink & \resultpending{APP_E3_SWITCH_N} &
    \resultpending{APP_E3_SWITCH_SUCCESS} &
    \resultpending{APP_E3_SWITCH_VALID} & \resultpending{APP_E3_SWITCH_COST} \\
    Cycle creation & \resultpending{APP_E4_CREATE_N} &
    \resultpending{APP_E4_CREATE_SUCCESS} &
    \resultpending{APP_E4_CREATE_VALID} & \resultpending{APP_E4_CREATE_COST} \\
    Cycle opening & \resultpending{APP_E4_OPEN_N} &
    \resultpending{APP_E4_OPEN_SUCCESS} &
    \resultpending{APP_E4_OPEN_VALID} & \resultpending{APP_E4_OPEN_COST} \\
    Fused/spiro/bridged reachable set & \resultpending{APP_E4_COMPLEX_N} &
    \resultpending{APP_E4_COMPLEX_SUCCESS} &
    \resultpending{APP_E4_COMPLEX_VALID} & \resultpending{APP_E4_COMPLEX_COST} \\
    \bottomrule
  \end{tabularx}
\end{table}

\subsection{E5--E6: quotient and exact-control verification}
\label{app:e5-e6-protocols}

\paragraph{E5 quotient interventions.}
For each source, the evaluator records every raw mark, executor successor,
canonical key, and aggregate mass.  The interventions are: a persistent-slot
permutation with action transport; a deterministic refinement of selected
marks into aliases whose masses sum to the original; and the inverse
coarsening.  We compare raw mark multiplicities, aggregate successor mass,
sampled molecular frequencies, and successor-level controlled rows.  A fresh
dictionary evaluator independently enumerates, executes, canonicalizes, and
groups marks.  Negative controls apply a mark-level power tilt and mark-level
top-$k$, which should change when alias count changes.

\paragraph{E6 graph.}
The \texttt{carbon\_6\_slots} graph contains 967 canonical states and 14,432
directed canonical edges and includes the distinguished null source.  It is a
bounded carbon-only verification language.  The graph audit requires
deterministic state indexing, exhaustive action disposition, explicit
vocabulary-boundary exclusions, no productive action outside the closed
slice, raw-mark-to-edge compression and alias counts, productive versus
virtual mass reconstruction, and production segmented aggregation against a
fresh dictionary oracle.  These checks establish the graph supplied to the
solver; they are not themselves an exact-control result.

The claim-bearing solver contract separately freezes the row-stochastic base
kernel, initial state or distribution, horizon $K$, terminal desirability
$g$, numeric tolerances, and a smaller explicit-path fixture.  Missing any one
of these fields blocks solver execution.  Exact values are obtained by
backward dynamic programming; $h_b(x)=0$ is represented exactly and the
controlled row is not numerically patched.  We compare the propagated
terminal law with the analytic tilted terminal law, verify the support subset
row by row, and restart the recursion from registered switch states for the
dynamic-continuation check.

\subsection{E7: same-base Pareto and dynamic control}
\label{app:e7-protocol}

\paragraph{Controller inputs.}
The learned value receives remaining budget, current molecule, immutable
source, goal, and protected/path-constraint context.  The base
editing prior remains frozen and goal independent.  Every same-base arm uses
the identical productive successor kernel and differs only in how supported
successors are scored, resampled, or backed up.

\paragraph{Pareto archive.}
All objectives are oriented so larger is better and normalized to $[0,1]$
using bounds frozen from training/development data or a fixed literature
convention.  The hypervolume reference is the origin of this normalized
maximization cube.  Raw objective values are retained.  The archive uses the
registered dominance convention, similarity floor, feasibility predicate, and
duplicate identity rule.  Primary endpoints are terminal normalized
hypervolume and hypervolume area under the oracle-call curve; IGD+, feasible
nondominated count, source similarity, structural diversity, and independent
held-out evaluator performance are secondary.

\paragraph{Dynamic switch.}
At the registered switch step, the new controller starts from the actual
current molecule and remaining budget.  Comparators continue under the new
goal, restart from the source, restart from the switch molecule under a
separate sampler, or pursue a static compromise.  Adaptation regret is
measured from the switch state at matched post-switch oracle cost.

\paragraph{Pareto fan.}
One valid prefix is branched toward multiple registered target regions.
Endpoint spread, region coverage, hypervolume, branch diversity, prefix reuse,
and total cost are compared with independent restarts.  Branches contribute
all feasible visited states to the common archive under one frozen
deduplication rule.

\paragraph{Path constraints.}
Scaffold, pharmacophore, atom-count, charge, and structural-alert constraints
are checked at every committed state.  Endpoint-only filtering is the matched
waste baseline.  A prefix-closed state constraint masks violating successors;
a history-dependent constraint augments controller state or kills the path.
We report path feasibility, forbidden intermediates, endpoint yield, and
wasted oracle calls.

\begin{table}[htbp]
  \centering
  \caption{\textbf{Dynamic-control diagnostics.} These complement the aggregate
  controller table in the main text.}
  \label{tab:appendix-dynamic}
  \small
  \begin{tabularx}{\linewidth}{@{}lYYYY@{}}
    \toprule
    Protocol & Success/coverage $\uparrow$ & Regret $\downarrow$ &
    Prefix reuse $\uparrow$ & Oracle cost $\downarrow$ \\
    \midrule
    Preference switch, learned value &
    \resultpending{APP_E7_SWITCH_SUCCESS} &
    \resultpending{APP_E7_SWITCH_REGRET} & n/a &
    \resultpending{APP_E7_SWITCH_COST} \\
    Restart from source &
    \resultpending{APP_E7_RESTART_SUCCESS} &
    \resultpending{APP_E7_RESTART_REGRET} & 0 &
    \resultpending{APP_E7_RESTART_COST} \\
    Pareto fan &
    \resultpending{APP_E7_FAN_COVERAGE} &
    \resultpending{APP_E7_FAN_REGRET} &
    \resultpending{APP_E7_FAN_REUSE} &
    \resultpending{APP_E7_FAN_COST} \\
    Pathwise control &
    \resultpending{APP_E7_PATH_SUCCESS} &
    \resultpending{APP_E7_PATH_WASTE} &
    \resultpending{APP_E7_PATH_PREFIX} &
    \resultpending{APP_E7_PATH_COST} \\
    Endpoint-only filter &
    \resultpending{APP_E7_FILTER_SUCCESS} &
    \resultpending{APP_E7_FILTER_WASTE} & n/a &
    \resultpending{APP_E7_FILTER_COST} \\
    \bottomrule
  \end{tabularx}
\end{table}

\subsection{E1: separate timed de novo protocol}
\label{app:e1-protocol}

The unconditional model uses a separate broad-organic corpus with nonzero de
novo mixture, a separate checkpoint, and the timed marked CTMC.  The source
distribution includes the null state and the reversible
$\varnothing\rightleftarrows$ one-atom interface.  Gillespie sampling uses the
calibrated productive molecular exit hazard; a fixed number of embedded edits
is not substituted for time.

Every model and baseline generates the registered number of molecules for
each of five seeds.  FCD uses the same ChemNet implementation and reference
sample count.  Held-out manifold recall uses frozen ChemNet embeddings and the
registered $k$-nearest-neighbor coverage construction.  Topology distance is
the equal-weight aggregate of graph cycle rank, ring-system count, ring-size
histogram, fused/spiro/bridged category, and aromatic-ring fraction distances.
All-step validity, endpoint validity, connectedness, and uniqueness are hard
gates; novelty, scaffold coverage, memorization, size/property distributions,
and throughput are secondary.

\subsection{Comparator definitions}
\label{app:comparators}

\begin{table}[htbp]
  \centering
  \caption{\textbf{Comparator roles.} Inclusion requires a runnable artifact
  and a compatible task contract; naming a method here does not imply that its
  published support equals COMPOSE support.}
  \label{tab:appendix-comparators}
  \small
  \begin{tabularx}{\linewidth}{@{}lYYY@{}}
    \toprule
    Comparator & Role & Support relation & Required accounting \\
    \midrule
    Uniform canonical successor &
    test learned transport &
    identical successor support &
    identical seeds, edit budget, executor \\
    Operator ablations &
    identify load-bearing birth/death/cycle families &
    preregistered support removal &
    otherwise identical prior/training or sampling contract \\
    Unguided, rerank, greedy, Boltzmann &
    local control baselines &
    identical frozen successor kernel &
    every scored successor and oracle call \\
    Scalarized and MOG-DFM-style guidance &
    preference-conditioned local guidance &
    identical frozen successor kernel &
    direction samples, forward passes, oracle calls \\
    SMC/Feynman--Kac &
    sequential population control &
    identical frozen successor kernel &
    particles, resampling, all oracle calls \\
    GraphGA, MARS &
    graph-search/valid-editor context &
    external, generally different &
    matched sources, objectives, constraints, budgets \\
    RetMol, GenMol &
    retrieval/discrete-generation context &
    external, generally different &
    independently runnable standardized endpoints \\
    \bottomrule
  \end{tabularx}
\end{table}

External rows are omitted when their released artifacts cannot honor the same
source panel, objective, similarity/constraint contract, and oracle budget.
Published numbers from a different standardization pipeline are contextual and
never enter paired tests.  One-shot methods receive ``not applicable'' for
pathwise metrics rather than an assumed zero violation rate.

\subsection{Statistics and reporting}
\label{app:statistics}

The independent resampling unit is the highest-level unit that can carry
shared chemistry: source/scaffold component for editing tasks and the
registered reference-sample component for distributional generation.
Trajectories, controller seeds, branches, and particles are nested replicates,
not independent molecules.  Primary method contrasts use paired
component-level bootstrap intervals.  Multiplicity correction is applied
within each preregistered family of primary comparisons; secondary diagnostics
are labeled exploratory.  We report the estimate, interval, number of
independent units, number of attempted trajectories, every seed-level result,
and all exclusions with reason codes.  Means across seeds are never presented
as corpus-level uncertainty.

\section{Reproducibility and audit inventory}
\label{app:reproducibility}

\subsection{Production primitives and experiment-only code}

Experiments reuse only the primitives that define the trained scientific
object: checkpoint reconstruction, legal enumerators, deterministic executor,
persistent-slot semantics, canonical molecular identity, operator/capability
configuration, and authoritative segmented successor aggregation.  Task
builders, exact graph construction, exact solver, controllers, learned-value
training, metrics, statistical analysis, and plots are written as experiment
code under the frozen registry.  Historical exact-Doob scripts, old exactness
JSON files, data-starved checkpoints, and pre-RingCore figures are not evidence.

\subsection{Required artifact chain}

\begin{table}[htbp]
  \centering
  \caption{\textbf{Claim-bearing artifact chain.} Every arrow is checked by
  physical SHA-256 and semantic identity, not inferred from filenames.}
  \label{tab:appendix-artifacts}
  \small
  \begin{tabularx}{\linewidth}{@{}lYY@{}}
    \toprule
    Stage & Required immutable artifacts & Final identity \\
    \midrule
    Corpus &
    source shards, compiler contract, packed shards, exclusion overlays,
    capability census, split audit &
    \resultpending{APP_ARTIFACT_CORPUS} \\
    Gate 0/T1 &
    operator freeze, exact panels, cache inventory, support/charge audit,
    threshold decision &
    \resultpending{APP_ARTIFACT_T1} \\
    Pilot/full training &
    recipe, ordered stream, monitor policy, recovery checkpoints, validation
    leaderboard &
    \resultpending{APP_ARTIFACT_TRAINING} \\
    Selected editing model &
    checkpoint bytes, reconstruction contract, capability and rollout report &
    \resultpending{APP_ARTIFACT_EDITING_MODEL} \\
    Exact control &
    graph audit, solver task, base kernel, values, controlled rows, path oracle &
    \resultpending{APP_ARTIFACT_E6} \\
    Controllers &
    frozen prior, tasks, objective normalization, oracle and compute logs &
    \resultpending{APP_ARTIFACT_E7} \\
    De novo &
    separate corpus, timed recipe, checkpoint, hazard calibration, samples &
    \resultpending{APP_ARTIFACT_E1} \\
    \bottomrule
  \end{tabularx}
\end{table}

\subsection{Known implementation boundaries}

The final release records four boundary groups explicitly.
\textbf{State semantics:} the null source is not a molecule, and persistent
padding is a coordinate representation rather than semantic fixed dimension.
\textbf{Chemical scope:} insertion supports zero or one existing neighbor;
formal charge is preserved rather than designed; and stereochemistry,
isotopes, radicals, explicit protonation changes, and 3D geometry are outside
the current state.  Chemical validity alone does not imply stability,
synthesis, potency, or safety.  \textbf{Topology and control:} cycle close/open
provides primitive topology support, while coordinated ring restatement is an
empirically retained accelerator; exact control is restricted to the bounded
enumerable graph, and molecular-scale value control is approximate.
\textbf{Regime separation:} editing and timed de novo checkpoints are not
interchangeable.
